\chapter{Đọc thêm 1: Đợt tăng vọt giá cổ phiếu GameStop dưới góc nhìn nhà đầu tư cá nhân (Zhao)}
\label{ch:M6L4DocThem1}

\section{Thông tin nguồn}

\begin{itemize}
  \item \textbf{Nguồn:} Zhao, Y. ``A Study of the GameStop Stock Price Surge During the Covid-19 Pandemic from the Perspective of Individual Investors.'' 2021.
  \item \textbf{Phần cần đọc:} Toàn bộ bài báo
  \item \textbf{Ghi chú:} bản chuyển ngữ tiếng Việt súc tích, bám cấu trúc nguồn (giữ nguyên tiêu đề, công thức, số liệu, bảng, chú thích và danh mục tài liệu tham khảo).
\end{itemize}

Nghiên cứu về đợt tăng vọt giá cổ phiếu GameStop trong đại dịch Covid-19 dưới góc nhìn của nhà đầu tư cá nhân

Yilin Zhao

Khoa Kinh tế, Đại học California, Los Angeles, 90024, Hoa Kỳ

*Tác giả liên hệ. E-mail: yilinzhao\_5025@yahoo.com

Đại dịch Covid-19 bùng phát tại Hoa Kỳ đã làm thị trường chứng khoán thay đổi đáng kể, đồng thời tỷ trọng giao dịch trực tiếp của nhà đầu tư cá nhân tăng lên. Vụ bán khống GameStop (GME) gần đây cho thấy sức mạnh của nhà đầu tư nhỏ lẻ, khiến người ta chú ý đến xu hướng quay lại của nhà đầu tư nhỏ lẻ trên thị trường cổ phiếu Mỹ. Bài báo nghiên cứu các yếu tố làm gia tăng hành vi giao dịch trực tiếp của nhà đầu tư cá nhân trên thị trường chứng khoán Mỹ trong thời kỳ Covid-19 và dự báo xu hướng chi phối của thị trường trong tương lai. Phương pháp nghiên cứu là phân tích dữ liệu từ Yahoo Finance, đồng thời tham chiếu các giả thuyết và chính sách kinh tế liên quan. Bốn yếu tố làm tăng giao dịch trực tiếp của nhà đầu tư cá nhân gồm: lệnh cách ly và các ưu đãi tài khóa sau khi dịch bùng phát, mô hình môi giới bán lẻ không thu phí hoa hồng, biến động giá cổ phiếu tạo thêm cơ hội đầu cơ, và ảnh hưởng của mạng xã hội. Nhà đầu tư cá nhân tạo ra tác động phức tạp lên thị trường chứng khoán; tuy nhiên, hành vi giao dịch trực tiếp tích cực của họ ít có khả năng kéo dài, và vai trò chi phối của nhà đầu tư tổ chức vẫn là xu hướng dài hạn.

Từ khóa: Nhà đầu tư cá nhân, Thị trường chứng khoán, Nền tảng giao dịch trực tuyến, Giao dịch trực tiếp, Giá tăng vọt

\section{1. GIỚI THIỆU}

Trong kỷ nguyên kinh tế số, thông tin lan truyền nhanh làm mở rộng phạm vi quan hệ quen biết và niềm tin, qua đó chia sẻ tiếng nói giữa các tổ chức chuyên nghiệp và nhà đầu tư cá nhân. Năm 2020, thị trường toàn cầu chứng kiến làn sóng giao dịch của nhà đầu tư cá nhân, với số tài khoản đăng ký mới kỷ lục tại bốn công ty môi giới bán lẻ lớn và tỷ trọng giao dịch tăng nhanh. Tháng 1/2021, cổ phiếu GameStop, một nhà bán lẻ trò chơi điện tử của Mỹ, bị nhà đầu tư cá nhân ép bán khống. Theo dữ liệu Yahoo Finance, giá cổ phiếu tăng vọt từ 18,84 USD cuối năm 2020 lên đỉnh 347,51 USD ngày 28/1/2021, tăng 134\% so với ngày hôm trước. Sau khi giảm về khoảng 40-50 USD vào giữa tháng 2 khi các nhà đầu tư chuyên nghiệp bắt đầu rút lui, giá cổ phiếu tiếp tục dao động mạnh. Thị trường vốn toàn cầu đang chú ý đến việc nhà đầu tư cá nhân tại Mỹ tích cực giao dịch cổ phiếu trực tiếp. Theo nghiên cứu của Ngoc {[}1{]}, năm yếu tố ảnh hưởng đến hành vi đầu tư của nhà đầu tư cá nhân gồm: hành vi bầy đàn, thị trường, triển vọng, tự tin thái quá - ngụy biện của con bạc, và thiên lệch neo giữ - năng lực; tuy nhiên có rất ít nghiên cứu giải thích vì sao nhà đầu tư cá nhân có thể đẩy giá một cổ phiếu lên mức cao như vậy. Từ góc nhìn của nhà đầu tư cá nhân, bài báo xem xét những yếu tố nào dẫn đến cuộc đối đầu giữa nhà đầu tư cá nhân và nhà đầu tư tổ chức xét về chính sách và công nghệ. Bài báo dùng dữ liệu từ Yahoo Finance và tham chiếu thông tin hiện có, kết hợp dữ liệu định lượng và định tính. Nghiên cứu này có thể giúp tối ưu hóa chiến lược đầu tư của tổ chức, và các yếu tố trên có thể dùng để dự báo lượng vốn đầu tư của nhà đầu tư cá nhân trong tương lai.

\section{2. BỐN YẾU TỐ GÂY RA ĐỢT TĂNG VỌT GIÁ CỔ PHIẾU GAMESTOP}

\subsection{2.1. Ưu đãi tài khóa}

Nhà đầu tư cá nhân, còn gọi là nhà đầu tư nhỏ lẻ, tự đầu tư vào chứng khoán và tài sản, thường với số tiền khiêm tốn. Cổ phiếu họ mua là một phần danh mục đầu tư cá nhân và không đại diện cho lợi ích của công ty nào. Nhiều nhà đầu tư cá nhân ra quyết định dựa trên cảm xúc, và một số yếu tố chi phối hành vi đầu tư của họ. Từ thập niên 1960-1970, với các quỹ tương hỗ và quỹ hưu trí đại diện cho nhà đầu tư tổ chức, nhà đầu tư tổ chức dần chiếm ưu thế trên thị trường chứng khoán Mỹ, nhưng nhà đầu tư cá nhân vẫn là những người nắm giữ đáng kể cổ phiếu Mỹ {[}2{]}. Sau khi đại dịch COVID-19 bùng phát do tính lây lan cao của virus, 42 bang và vùng lãnh thổ đã ban hành lệnh ở nhà từ 19/3/2020 đến 31/5/2020, ảnh hưởng đến 2.355 (73\%) trong số 3.233 quận của Mỹ {[}3{]}. Lệnh cách ly tại nhà hạn chế phạm vi sinh hoạt bình thường của người dân. Ba đợt thanh toán hỗ trợ kinh tế là khoản tiền do IRS phát hành để giúp người dân trong đại dịch, cùng với các chính sách giảm thuế khác liên quan đến virus corona. Theo số liệu của Cục Phân tích Kinh tế Hoa Kỳ, trong tháng 4/2020, thu nhập cá nhân tăng 1,97 nghìn tỷ USD (10,5\%), thu nhập khả dụng cá nhân (DPI) tăng 2,13 nghìn tỷ USD (12,9\%), trong khi chi tiêu tiêu dùng cá nhân (PCE) giảm 1,89 nghìn tỷ USD (13,6\%) {[}4{]}. Số liệu cho thấy dịch bệnh làm ngành dịch vụ ngoại tuyến co hẹp và hạn chế

các hành vi tiêu dùng hằng ngày như ăn uống và giải trí, những cư dân bị cách ly tại nhà có xu hướng rót nhiều vốn hơn vào thị trường chứng khoán vì có nhiều thời gian rảnh hơn so với cuộc sống bình thường. Đại dịch cũng khiến nhiều người thất nghiệp: tháng 4/2020 tỷ lệ thất nghiệp là 14,8\%, mức cao nhất kể từ khi bắt đầu thu thập dữ liệu năm 1948. Đến tháng 5/2021, tỷ lệ thất nghiệp vẫn cao hơn (5,8\%) so với tháng 2/2020 (3,5\%){[}5{]}. Centiment Research khảo sát 2.054 cá nhân từ ngày 24 đến 26/6/2020 và thấy 83\% người được hỏi nhận séc của chính phủ, trong đó một nửa đưa một phần tiền hỗ trợ vào thị trường chứng khoán. Người thất nghiệp có động lực mạnh để kiếm tiền, và giao dịch chứng khoán trực tuyến rõ ràng rất thuận tiện. Sau khi COVID-19 bùng phát, lệnh ở nhà và hỗ trợ tài chính đã thúc đẩy nhà đầu tư cá nhân mở tài khoản và tham gia thị trường.

Figure 1: Personal Income and Outlays: April 2020 (Thu nhập cá nhân và chi tiêu: tháng 4/2020)

\subsection{2.2. Mô hình môi giới bán lẻ không hoa hồng}

Các nhà môi giới bán lẻ thường hỗ trợ khách hàng giao dịch cổ phiếu và nợ, như cổ phiếu và trái phiếu. Từ tháng 9/2019, một số công ty môi giới lớn đã hạ hoa hồng về 0 cho các giao dịch, khiến việc giao dịch của nhà đầu tư cá nhân cực kỳ rẻ và giúp các công ty phát triển các mảng kinh doanh phụ trợ. Theo ước tính của Bloomberg Intelligence, nhà đầu tư cá nhân sẽ chiếm 20\% tổng số cổ phiếu được giao dịch tại Hoa Kỳ năm 2020; sau quý I/2021, con số này đã tăng lên 22\%. Việc chuyển sang giao dịch không hoa hồng là lợi thế cho nhà đầu tư, đặc biệt là nhà đầu tư mới có vốn hạn chế. Các nghiên cứu cho thấy các nhà môi giới không hoa hồng làm tăng đáng kể thị phần tài sản của khách hàng, và việc gia tăng các nhóm lệnh quy mô nhỏ trong khối lượng từ lệnh bán lẻ có liên quan đến việc không thu hoa hồng{[}6{]}. Khi các nhà môi giới bán lẻ hạ hoa hồng, nhà đầu tư cá nhân sẵn sàng tham gia thị trường hơn. Khi ngày càng nhiều nhà đầu tư và nhà giao dịch dùng máy tính bảng hoặc điện thoại thông minh làm nền tảng chính, các công ty này đã cải thiện ứng dụng di động để thu hút thêm người dùng di động. Các ứng dụng đầu tư trực tuyến tốt nhất mang lại trải nghiệm thống nhất giữa máy tính để bàn và thiết bị di động, gồm khả năng chia sẻ danh sách theo dõi và cảnh báo, bộ lọc cổ phiếu và nộp séc qua ảnh. Chẳng hạn, nền tảng môi giới trực tuyến Robinhood, nổi tiếng với giao dịch miễn hoa hồng tiên phong, ra mắt năm 2015 qua ứng dụng di động, đã tăng số người dùng từ 0,5 triệu năm 2015 lên 18 triệu năm 2021{[}7{]}. Nhờ giao dịch trực tuyến không hoa hồng và sự tiện lợi của phần mềm, nhà đầu tư cá nhân có thể dễ dàng đầu tư trực tiếp vào thị trường chứng khoán. Ngoài ra, xu hướng cho phép khách hàng mua bán cổ phiếu lẻ (fractional shares) thu hút nhiều nhà đầu tư trẻ và giúp họ đa dạng hóa danh mục đầu tư với vốn tương đối nhỏ.

Figure 2: Robinhood users between 2015 to 2021 (Số người dùng Robinhood giai đoạn 2015-2021)

Figure 3: Robinhood revenue between 2015 to 2020 (Doanh thu Robinhood giai đoạn 2015-2020)

\subsection{2.3. Biến động cổ phiếu làm gia tăng cơ hội đầu cơ trên thị trường}

Biến động (volatility) là tốc độ giá cổ phiếu tăng hoặc giảm trong một khoảng thời gian nhất định. Biến động giá cổ phiếu càng lớn thường kéo theo rủi ro càng cao, đồng thời giúp nhà đầu tư dự đoán các dao động trong tương lai.

S\&P 500 (Standard and Poor\textquotesingle s 500) là chỉ số cổ phiếu theo dõi hiệu quả hoạt động của 500 công ty lớn, được dùng làm ví dụ để nghiên cứu biến động của thị trường chứng khoán Hoa Kỳ và là một trong những chỉ số cổ phiếu được theo dõi phổ biến nhất. Chỉ số này sẽ được dùng để xem xét biến động cổ phiếu sau khi COVID-19 bùng phát. Biến động của thị trường chứng khoán cũng có thể xuất hiện do thông tin mới, bất ngờ làm thay đổi lợi suất cổ phiếu kỳ vọng {[}8{]}.

Có bằng chứng cho thấy lợi suất của Dow Jones và S\&P có tác động thuận chiều lên phương sai có điều kiện thay đổi do COVID-19 {[}9{]}. Biến động thị trường mang lại cho nhà đầu tư cơ hội sinh lời; mua ở đáy không mất nhiều thời gian để thu lợi. Trong giai đoạn biến động cao như đại dịch COVID-19, nhà đầu tư dễ ra quyết định thiên lệch hoặc cảm tính hơn. Biến động của thị trường chứng khoán Hoa Kỳ năm 2020 thu hút nhà đầu tư cá nhân tham gia để tìm kiếm lợi suất vượt trội. Dù nhiều nhà đầu tư cá nhân quen đầu tư vào quỹ tương hỗ và quỹ chỉ số hoán đổi danh mục (ETF) để tránh biến động, nhưng do đại dịch, sự phân hóa tình hình kinh doanh của các công ty niêm yết tại Hoa Kỳ gây ra dao động giá cổ phiếu, thu hút nhà đầu tư cá nhân tìm kiếm cơ hội giao dịch đầu cơ. Ngoài ra, xét theo phân bố giá trị của thị trường, trong các cổ phiếu cấu thành S\&P 500, một số gã khổng lồ công nghệ chiếm gần một phần tư tổng giá trị vốn hóa, nên nhà đầu tư cá nhân chỉ cần đầu tư vào các cổ phiếu dẫn đầu cũng có thể đạt lợi suất đáng kể tương tự chỉ số. Nhiều nhà đầu tư cá nhân có thể sẵn sàng phân bổ vào cổ phiếu riêng lẻ hơn, qua đó thúc đẩy thêm thị trường bán lẻ.

2.4. Ảnh hưởng của mạng xã hội

Nhà đầu tư cá nhân ngày càng dựa vào lời khuyên đầu tư trên các nền tảng mạng xã hội, ngay cả khi những lời khuyên đó có rất ít giá trị dự báo vẫn có thể tác động đến quyết định của họ. Khi cổ phiếu xuất hiện trên tin tức, nhà đầu tư thực bị truyền thông chi phối và có xu hướng mua hơn là bán, nên việc mua dựa trên sự chú ý này có thể khiến họ giao dịch quá mức mang tính đầu cơ, ảnh hưởng đến giá cổ phiếu {[}10{]}. Các nhà nghiên cứu phát hiện một số nhà đầu tư dựa vào lời khuyên có giá trị dự báo thấp vì tin rằng nó mang tính chỉ dẫn, trong khi số khác không nhận ra tác động của lời khuyên đó lên quyết định đầu tư của mình {[}11{]}. Việc kêu gọi trên mạng xã hội có thể dẫn đến hành vi bầy đàn (herding). Hành vi bầy đàn xảy ra khi một nhóm nhà đầu tư có cùng khuynh hướng làm theo hành vi giao dịch của một người dẫn dắt {[}12{]}.

Nhờ tiến bộ công nghệ, mạng xã hội thường tràn ngập thông tin đầu tư, có thể dẫn dắt nhà đầu tư cá nhân mua một cổ phiếu nhất định. Một nhân tố đằng sau đợt tăng vọt giá cổ phiếu GameStop là ảnh hưởng của mạng xã hội, khi nhiều người kêu gọi cùng mua cổ phiếu để tránh việc công ty sụp đổ. Angel cho rằng, theo nghiên cứu, một số nhà đầu tư chia sẻ lợi suất cổ phiếu và gợi ý đầu tư qua mạng xã hội, đóng vai trò kích động đối với các nhà đầu tư cá nhân thiếu nguồn thông tin và khả năng phân tích, phân biệt {[}13{]}.

Cá nhân tham gia hành vi bầy đàn khi họ làm theo người khác thay vì đưa ra phán đoán dựa trên chuyên môn và dữ kiện của chính mình. Hành vi bầy đàn có thể dẫn đến những hành động phi lý ở các nhóm lớn. Trong các bong bóng và đợt sụp đổ thị trường, đây là hiện tượng bầy đàn thường gặp ở nhà đầu tư. Sự lạc quan phi lý đẩy giá tài sản vượt xa giá trị cơ bản trong thời kỳ bong bóng {[}14{]}. Cảm xúc (affect) là lối tắt trong ra quyết định, trong đó tình cảm tích cực hoặc tiêu cực của một người tác động đến lựa chọn của họ. Đây là phản ứng cảm xúc trước thông tin chứ không phải phản ứng lý tính. Cảm xúc có xu hướng làm thay đổi nhận thức của con người về lợi ích và rủi ro của các quyết định. Nói chung, nếu thông tin về một quyết định khiến ta hài lòng, ta có nhiều khả năng đánh giá cao lợi ích của lựa chọn đó. Chẳng hạn, trên diễn đàn Reddit WallStreetBets, một nhà đầu tư cá nhân đi đầu đã đăng câu chuyện làm giàu nhanh chóng, cho thấy khoản đầu tư ban đầu 4.500 USD thành 1,5 triệu USD, nhằm thu hút "đám đông bán lẻ". Do đó, thông tin chứng khoán và đầu tư trên mạng xã hội là một lý do nữa khiến nhà đầu tư cá nhân tham gia giao dịch trực tiếp trên thị trường chứng khoán.

\begin{enumerate}
\def\labelenumi{\arabic{enumi}.}
\setcounter{enumi}{2}
\item
  THẢO LUẬN
\end{enumerate}

Về dài hạn, nhiệt huyết của nhà đầu tư cá nhân sẽ không nhanh chóng phai nhạt, nhưng vẫn còn khoảng cách so với vị thế thống trị trên thị trường đầu tư. Người ta cho rằng tỷ trọng giao dịch của nhà đầu tư cá nhân ở Hoa Kỳ dự kiến tiếp tục tăng, song nhà đầu tư tổ chức vẫn sẽ chi phối thị trường, vì các lý do sau. Thứ nhất, thói quen tiêu dùng và đầu tư hình thành trong đại dịch có thể tiếp tục, nên một số nhà đầu tư cá nhân Mỹ sẽ không rút khỏi thị trường chứng khoán ngay sau đại dịch. Thứ hai, xu hướng môi giới không hoa hồng đã hình thành, và mô hình giao dịch không hoa hồng trong tương lai sẽ tiếp tục nuôi dưỡng nhiệt huyết của nhà đầu tư bán lẻ. Thứ ba, giá trị vốn hóa của các công ty hàng đầu Hoa Kỳ vẫn đang mở rộng, và ảnh hưởng của họ lên chỉ số cũng tăng. Nếu thị trường chứng khoán Mỹ vẫn mạnh, nhà đầu tư cá nhân có thể sẵn sàng đầu tư vào cổ phiếu riêng lẻ hơn là quỹ.

\begin{enumerate}
\def\labelenumi{\arabic{enumi}.}
\setcounter{enumi}{3}
\item
  KẾT LUẬN
\end{enumerate}

Từ góc nhìn của nhà đầu tư cá nhân, bài báo này xem xét các nguyên nhân có thể khiến giá cổ phiếu GameStop tăng, gồm: cách ly và các biện pháp kích thích tài khóa sau khi dịch bùng phát, mô hình môi giới bán lẻ không hoa hồng cùng sự trỗi dậy của các nền tảng giao dịch trực tuyến, biến động cổ phiếu làm tăng cơ hội đầu cơ trên thị trường, và ảnh hưởng của mạng xã hội. Bốn yếu tố bên ngoài này đều có tác động nhất định đến việc nhà đầu tư cá nhân mua cổ phiếu quy mô lớn, dẫn đến giá cổ phiếu tăng mạnh. Sự phát triển của fintech giúp cá nhân tham gia trực tiếp vào thị trường chứng khoán dễ hơn, và đại dịch COVID-19 đã thay đổi phần lớn cách chúng ta nhìn nhận mọi việc. Về dài hạn, nhà đầu tư tổ chức vẫn sẽ chi phối thị trường, còn nhà đầu tư cá nhân có thể tham gia thị trường khi xu hướng lợi nhuận thị trường chưa rõ ràng, như việc nhà đầu tư cá nhân quy mô lớn rút lui do khủng hoảng tài chính 2008. Trong kỷ nguyên hậu dịch, ảnh hưởng chồng chéo và lên men của hai yếu tố này đến hành vi đầu tư của nhà đầu tư cá nhân vẫn cần nhiều quan sát và nghiên cứu. Hạn chế của nghiên cứu là xét từ góc độ nhà đầu tư không phải cá nhân, giá cổ phiếu tăng nhanh cũng có thể là kết quả của việc nhà đầu tư tổ chức thao túng nhà đầu tư cá nhân, điều này chưa được xem xét trong bài báo. Bài báo không đưa ra phân tích định lượng về tác động của các yếu tố này đến hành vi tiêu dùng, và các nghiên cứu tiếp theo có thể tập trung vào phân tích định lượng tác động của các yếu tố bên ngoài đến hành vi đầu tư cá nhân.

REFERENCE

{[}1{]} Ngoc, Luu Thi Bich. "Behavior pattern of individual investors in stock market." International Journal of Business and Management 9.1 (2014): 1.

{[}2{]} Davis E P, Steil B. Institutional investors{[}M{]}. Massachusetts: The MIT Press, 2004.

{[}3{]} Moreland A, Herlihy C, Tynan MA, et al. Timing of State and Territorial COVID-19 Stay-at-Home Orders and Changes in Population Movement --- United States, March 1--May 31, 2020. MMWR Morb Mortal Wkly Rep 2020;69:1198--1203. DOI: http://dx.doi.org/10.15585/mmwr.mm6935a2

{[}4{]} "Personal Income and Outlays: April 2020" Bureau of Economics Analysis https://www.bea.gov/news/2020/personal-incomeand-outlays-april-2020.

{[}5{]} Falk, Romero., etal. " Unemployment Rate during Covid-19 Pandemic." Congressional Research Service. Updated June 15, 2021 https://fas.org/sgp/crs/misc/R46554.pdf.

{[}6{]} Jain, Pankaj K. and Mishra, Suchismita and O\textquotesingle Donoghue, Shawn and Zhao, Le, Trading Volume Shares and Market Quality in a Zero Commission World (December 2, 2020). Available at SSRN: https://ssrn.com/abstract=3741470 or http://dx.doi.org/10.2139/ssrn.3741470.

{[}7{]} Curry, David. "Robinhood Revenue and Usage Statistics (2021)" July 11, 2021. Business of Apps. https://www.businessofapps.com/data/robinhoodstatistics/.

{[}8{]} Engle, R.F., Victor, K.Ng. (1993), Measuring and Testing the impact on Volatility. The Journal of Finance, 48(5), 1749-1778.

{[}9{]} Onali, Enrico. "COVID-19 and stock market volatility." Available at SSRN 3571453 (2020).

{[}10{]} Barber, Brad M., and Terrance Odean. "The behavior of individual investors." Handbook of the Economics of Finance. Vol. 2. Elsevier, 2013. 1533-1570.

{[}11{]} Kadous, Kathryn and Mercer, Molly and Zhou, Yuepin, Do Individual Investors Understand How Social Media Advice Influences Their Investment Decisions? (June 17, 2019). Available at SSRN: https://ssrn.com/abstract=2968407 or http://dx.doi.org/10.2139/ssrn.2968407.

{[}12{]} Shu-Fan Hsieh, Chia-Ying Chan, Ming-Chun Wang, Retail investor attention and herding behavior, Journal of Empirical Finance, Volume 59, 2020, Pages 109-132, ISSN 0927-5398, https://doi.org/10.1016/j.jempfin.2020.09.005.

{[}13{]} Angel J J. Gamestonk: what happened and what to do about it{[}R{]}. Georgetown University Working Paper, 2021.

{[}14{]} Dale, Richard S., et al. "Financial Markets Can Go Mad: Evidence of Irrational Behaviour during the South Sea Bubble." The Economic History Review, vol. 58, no. 2, 2005, pp. 233--271. JSTOR, www.jstor.org/stable/3698692. Accessed 27 Aug. 2021.
